% --------------------------------------------------------------- thank you
% the thanks where the title slide has its title; drawn inside the frame, so that the notes' thumbnail
% shows it too
\begin{frame}[plain,noframenumbering]
\vbox to 0pt{\vskip-2mm\hbox to 0pt{\hskip-8mm\begin{tikzpicture}[x=1mm,y=1mm]
  \useasboundingbox (0,0) rectangle (160,-90);
  \node[anchor=north west,inner sep=0pt] at (8,-8) {\includegraphics[width=30mm]{\udslogo}};
  \node[anchor=north west,inner sep=0pt,text width=88mm,align=left] at (8,-36) {%
    {\fontsize{32}{36}\selectfont\bfseries\color{udsblue}Thank you\par}\vskip4mm
    \textcolor{udsred}{\rule{14mm}{0.8pt}}\par\vskip3mm
    {\fontsize{11}{13}\selectfont\color{udsdark}\insertauthor\par}};
\end{tikzpicture}\hss}\vss}
\par\rule{0pt}{80mm}
\note{\textbf{23:40--24:00}\par
    Thank you.\par
    Questions to open the discussion: would edge and SSIM terms help our segmentation of grain
      boundaries; what minimum grain size do we use in EBSD analysis, and why; how many independent
      specimens would we want before trusting a property model; and what would convince us that a
      generated microstructure is new and can be made.\par
    Backup slides follow: the loss terms in full, the 25 processing routes and the references.}
\end{frame}
